\chapter{Introduction}

\section{Background of the Study}
Robotics engineering has transitioned rapidly from isolated industrial work-cells toward dynamic, human-centered environments where machines and humans share physical workspaces. In these collaborative settings, autonomous mobile robots are required to operate intelligently alongside human personnel rather than merely executing rigid, pre-programmed trajectories in segregated areas. This operational shift has established Human-Robot Interaction (HRI) as a foundational discipline in modern robotics, requiring machines to interpret human social cues and intentions with high operational reliability. For an autonomous mobile robot to function safely and intuitively in shared spaces, it must recognize, interpret, and respond to human behavior in a manner that is predictable, natural, and immediate.

A primary communication channel in human-to-human collaboration is the use of natural hand gestures and body movement cues. Although modern computer vision algorithms have made optical gesture recognition increasingly practical, conventional deep learning implementations typically rely on cloud infrastructure or power-intensive Graphics Processing Units (GPUs) \cite{mahmud2022}. This reliance introduces severe communication latencies and vulnerability to network interruptions, rendering remote processing impractical for time-critical robotic safety applications \cite{tsitos2022}. Furthermore, a mobile robotic system cannot operate effectively by merely classifying gestures in isolation; it must maintain active spatial awareness of its surrounding environment to execute navigational directives safely. This study addresses these interconnected challenges by engineering and physically deploying a unified, vision-based intention recognition and autonomous navigation pipeline operating strictly on an embedded edge computing platform.

\section{Problem Statement}
Despite significant advancements in collaborative robotics, achieving seamless, responsive, and robust human-robot collaboration on low-cost, resource-constrained hardware remains an unresolved engineering bottleneck. Most industrial and commercial service robots still depend on rigid control modalities, such as physical teach pendants, capacitive touchscreens, or acoustic voice commands. These conventional interfaces frequently prove ineffective in noisy industrial manufacturing facilities or sterile healthcare environments where physical contact must be strictly avoided.

A critical analysis of state-of-the-art vision architectures reveals several practical barriers that impede reliable field deployment. Although laboratory prototypes frequently demonstrate high benchmark accuracy, their real-world reliability deteriorates significantly under unconstrained operational conditions. Specifically, existing vision-based intention recognition frameworks exhibit three critical technological limitations when transitioned to real-world embedded deployment:
\begin{enumerate}
    \item \textbf{High-Power Hardware and Cloud Dependency:} Many contemporary vision-based gesture architectures rely heavily on high-end desktop GPUs or offboard cloud computing services for inference \cite{mahmud2022}. In mobile robotics, reliance on offboard processing introduces unpredictable network latency, consumes excessive power, and introduces single-point network failure risks that violate real-time safety thresholds \cite{tsitos2022}.
    \item \textbf{Spatial Coordinate Domain Shift Across Distances:} Standard machine learning classifiers frequently evaluate raw anatomical pixel coordinates extracted from image frames. Because absolute landmark coordinates vary drastically as an operator moves closer to or farther from the camera, conventional models suffer severe classification degradation under real-world operational distance variations.
    \item \textbf{Ungrounded and Spatial-Blind Command Execution:} Existing gesture-controlled mobile robots predominantly execute motor commands blindly without contextual awareness of their physical surroundings. Translating recognized gestures into motor velocities without continuous laser-based environment mapping and obstacle verification presents immediate collision hazards in populated indoor workspaces.
\end{enumerate}

To overcome these engineering limitations, there is a distinct requirement for a lightweight, edge-deployable framework that tightly couples real-time intention recognition with spatial mapping within a unified, standardized robotic middleware. Such a framework must operate deterministically within embedded compute and thermal constraints while maintaining uncompromised collaborative safety. Furthermore, bridging the gap between high-level visual perception and low-level chassis actuation is essential to prevent latency bottlenecks that could lead to erratic or dangerous vehicular maneuvers.

\section{Objectives of the Study}

\subsection{General Objective}
The general objective of this research is to design, implement, and empirically evaluate an edge-deployed, real-time human intention recognition and navigation system that integrates scale-invariant hand gesture classification, temporal motion prediction, and spatial mapping into a unified ROS2 middleware pipeline on an autonomous mobile robot. This involves establishing a complete perception-to-action control loop that executes reliably on a physical mobile testbed without dependence on remote servers. In doing so, the study demonstrates the feasibility of deterministic human-robot collaboration under realistic indoor operational constraints.

\subsection{Specific Objectives}
Translating the broad engineering goal into a functioning cyber-physical system requires a structured multi-stage development methodology. Each phase addresses a specific perception, decision-making, or navigation challenge within the embedded compute envelope. To achieve this general objective, the following specific objectives were established:
\begin{enumerate}
    \item To design and implement an edge-optimized computer vision pipeline utilizing spatial zone filtering, anatomical hand landmark tracking, and a Multilayer Perceptron (MLP) classifier trained on nineteen engineered, scale-invariant geometric hand features.
    \item To construct, balance, and validate a custom 6,000-sample gesture dataset captured directly through the onboard monocular camera to eliminate sensor domain shift across six distinct operational commands.
    \item To develop a sequential motion prediction model utilizing a Long Short-Term Memory (LSTM) recurrent neural network to forecast short-horizon operator movement trajectories from continuous body pose landmarks.
    \item To implement a centralized, deterministic decision-making node (\texttt{brain\_node}) featuring temporal majority-vote confirmation, active subject-locking, and proportional visual servoing within a multi-node ROS2 Humble software architecture.
    \item To configure and physically validate Simultaneous Localization and Mapping (SLAM) and autonomous navigation stacks (Nav2) on the mobile platform to support safe, spatially aware trajectory execution.
    \item To quantitatively evaluate system performance through classification accuracy, confusion matrix analysis, component latency budgets, and real-time hardware throughput measurements on an embedded edge processor.
\end{enumerate}

\subsection{Engineering Research Questions}
To provide rigorous scientific evaluation, the experimental investigation is guided by formal engineering inquiries that test core hypotheses regarding system performance. These inquiries systematically interrogate the mathematical boundaries of the feature representation, computational latency, and collaborative safety. In alignment with Faculty of Engineering standards, this study formulates four core Engineering Research Questions (ERQs) to evaluate the technical feasibility, mathematical invariance, multi-user robustness, and safety compliance of the deployed system:
\begin{enumerate}
    \item \textbf{ERQ 1 (Edge Feasibility and Latency Budget):} Can a multi-stage visual perception pipeline—incorporating person detection, twenty-one-point hand landmark extraction, and geometric feature computation—execute natively on a low-cost, low-power edge processor (Raspberry Pi 5) without GPU acceleration while maintaining an end-to-end processing cycle below the 150~ms human-interaction latency threshold?
    \item \textbf{ERQ 2 (Mathematical Feature Invariance):} To what extent does converting raw anatomical landmark coordinates into scale- and distance-invariant geometric features (joint curl angles and palm-normalized Euclidean distances) eliminate spatial distortion and improve gesture classification accuracy across varying operator distances compared to unnormalized coordinate baselines?
    \item \textbf{ERQ 3 (Multi-Person Workspace Disambiguation):} How effectively can a central geometric spatial-zone filter coupled with rolling majority-vote temporal buffering isolate and maintain an active control lock on a primary operator in a shared workspace, preventing unintended command execution from background bystanders?
    \item \textbf{ERQ 4 (Intention-to-Action Coupling and Safety Verification):} Can discrete hand gesture commands and continuous kinematic motion trajectories be reliably integrated within ROS2 middleware to govern mobile robot velocity commands in full compliance with ISO 15066 collaborative robotic safety and reaction requirements?
\end{enumerate}

\section{Scope of the Study}
The operational scope of this research is bounded to the design, software integration, and physical deployment of a gesture-directed autonomous mobile robot operating within indoor environments. The testing environment is restricted to structured indoor settings with stable floor surfaces and standard ambient fluorescent illumination to ensure predictable sensor characteristics. To maintain deterministic control and prevent command ambiguity, the visual interaction vocabulary is restricted to six predefined static hand gestures (STOP, GO, LEFT, RIGHT, BACK, FOLLOW) and five discrete body motion states (Approaching, Moving Away, Moving Left, Moving Right, Stationary). 

Interaction is programmatically constrained to a single validated human operator at any given time, enforced through spatial-zone acceptance filtering and temporal lock scheduling. Furthermore, to address network vulnerability and communication latency constraints, the entire perception, intention interpretation, and motor control software stack operates strictly onboard the robot's local processor without any reliance on remote cloud servers or offboard computing clusters. Real-time operation is defined as the capability of the onboard architecture to sustain an end-to-end sensing-to-actuation cycle sufficient to maintain responsive, jitter-free interaction at human operating speeds.

\section{Significance of the Study}
The engineering significance of this study lies in demonstrating that complex, multi-modal human intention recognition can be successfully decoupled from power-hungry workstation GPUs and external cloud infrastructure. By deploying lightweight, mathematically engineered feature extractors directly onto embedded edge hardware, this research provides a scalable, cost-effective blueprint for autonomous mobile robots operating in resource-constrained settings. This architectural approach dramatically reduces deployment expenses and operational power consumption while eliminating vulnerability to wireless network dropouts.

From a practical perspective, an autonomous mobile robot capable of interpreting touchless gestures and body motion addresses critical requirements across several industrial and civil domains. In acoustically noisy manufacturing environments, optical gesture recognition provides an intuitive, reliable control alternative where acoustic voice recognition routinely fails. In sterile medical facilities and cleanrooms, touchless optical interaction prevents physical contamination and preserves strict hygiene protocols by eliminating contact with shared touchscreens. Academically, this study contributes an open, modular ROS2 architecture that combines explicit gesture classification with predictive motion tracking and laser-based spatial mapping. Ultimately, this work supports safer, more intuitive human-robot collaboration by ensuring predictable machine responses that adhere to international safety standards for collaborative robotic workspaces \cite{iso150662016}.

\section{Organization of the Study}
This project report is systematically structured into five cohesive chapters in accordance with the university project guidelines. Chapter One introduces the research background, defines the specific engineering bottlenecks, formulates the general and specific objectives alongside the core Engineering Research Questions, and delineates the project's scope and significance. Chapter Two presents a comprehensive review of related literature, examining theoretical foundations in Human-Robot Interaction, analyzing prior works in gesture classification and intention-aware navigation, and identifying the precise technological gaps addressed by this research. 

Chapter Three details the complete system design and methodology, encompassing physical hardware specifications, Docker-based containerized software deployment, geometric feature mathematics, LSTM path predictor training, and centralized state machine architecture. Chapter Four presents the empirical testing procedures, experimental results, classification confusion matrices, latency budgets, hardware throughput benchmarks, and comparative literature benchmarking. Finally, Chapter Five concludes the thesis by synthesizing the primary engineering findings, discussing physical implementation constraints, and proposing concrete recommendations for future industrial extensions.
